\talkslide{attacker-objective}{Learning the attacker: online DPO}{Timing: provisional}{Judge the outcomes of two attacks to construct a preference pair for the attacker.}{
\preferenceflow{0}

}
\speak{The judge compares defender responses, but we can use the result to prefer one attack over another. Start from the same seed and generate two candidate attacks. Hold the defender parameters fixed and examine its responses. The two defender boxes are independent calls to the same model, not separate models or a shared bottleneck. If attack $x_2$ induces the more successful attack outcome, we prefer $x_2$ to $x_1$ for the attacker. The seed is $x_s$, and $y_1,y_2$ are the corresponding sampled defender responses. The example shows an illustrative winner, not a measured attack. On harmful seeds the attacker seeks harmful compliance; on benign seeds it seeks unnecessary deflection. Once we have preferred and rejected attacks, online DPO can train the attacker to produce more of the preferred ones. This is the same preference-learning idea as for responses, with attacks as the objects being optimized. For clarity the diagram shows one response per attack; the full AdvGame construction first checks faithfulness and compares each attack through its preferred defender response. The defense section will show that construction explicitly. Source: AdvGame, Section 3.2 and Algorithm 1.}
